%! Author = taj
%! Date = 7/22/25

% Preamble
\documentclass[11pt]{article}

% Packages
\usepackage{amsmath}
\usepackage{graphicx}
\usepackage{float}

\title{6. Merjenje specifične toplote vode}
\author{Taj Šobot}

% Document
\begin{document}
    \maketitle
    \flushleft
    Pri tem eksperimentu smo določali specifične in toplotne kapacitete vode in predmetov.


    \section{Meritve in rezultati}
    V kalorimeter smo nalili vodo z začetno temperaturo $60 ^\circ$.
    Potem pa smo počakali, da se temperatura vode in kalorimetra v ravnovesju.

    Sobna temperatura oziroma začetna temperatura kalorimetra:
    \[
        T_s = T_k^{(z)} = (296 \pm 1) K
    \]

    Začetna temperatura vode:
    \[
        T_v^{(z)} = (333 \pm 1) K
    \]

    Ko smo vodo dodali v kalorimeter in počakali, da se temperatura izenači (28 min 20 s), je sistem bil na končni temperaturi
    \[
        T^{(k)} = (331 \pm 1) K
    \]

    Potem smo sistem kalorimeter + voda začeli segrevat.
    Grelec je deloval z tokom $I = (3.0 \pm 0.2) A$ in na napetosti $U = (20.1 \pm 0.1) V$.
    Grelec je deloval za $t = (1519 \pm 1)s $.
    Po segrevanjem končno temperaturo.
    \[
        T^{(kk)} = (346 \pm 1 )K
    \]
    Masa vode je bila izmerjena:
    \[
        m_v = (1,152 \pm 0,002)kg
    \]
    Te podatke lahko vstavimo v enačbo (8) iz skripte:
    \[
        c = \frac{I U t}{m_v \Delta T(t)\left(1+ \frac{T_v^{(z)} - T^{(k)}  }{ T^{(k)}  - T_k^{(z)}}\right)} = (5100 \pm 400) \frac{J}{kg K}
    \]
    To pa lahko vstavimo v enačbo (7):
    \[
        C_k = m_v c \frac{ T_v^{(z)} - T^{(k)}}{T^{(k)} - T_k^{(z)}} = (250 \pm 90) J/K
    \]

    \section{Vprašanja in odgovori}
    \textit{1. Zakaj lahko v enačbah (2) in (3) vstavimo delta T v kelvinih ali stopinjah celzija?}
    \\
    Razlika med kelvini in stopinjami je samo premaknjena skala.
    Razlike ($\Delta T$) so zato neodvisne od enote.
    \\
    \textit{2. Specifična toplota snovi A je večja od specifične toplote snovi B. Obe snovi imata enako
    začetno temperaturo, nato pa jima dovedemo enako količino toplote. Predpostavimo, da nobena
    od snovi ne spremeni agregatnega stanja.}
    \\
    V primeru $T_A < T_B$.
    "A" ima večjo specifično toploto, to pomeni da se bo za enoto dovedene energije pridobila manj temperature.
    \\
    \textit{3. Specifična toplota ledu je približno enaka polovici specifične toplote vode. Enaki masi ledu
    in vode dovedemo enako količino toplote. Začetna temperatura ledu je dovolj nizka, da se led
    med poskusom ne tali. Določite razmerje med spremembama temperature vode in ledu.}
    \\
    Temperatura ledu bo narastla 2x hitreje.

\end{document}